% Options for packages loaded elsewhere
\PassOptionsToPackage{unicode}{hyperref}
\PassOptionsToPackage{hyphens}{url}
\documentclass[
  ignorenonframetext,
  aspectratio=169,
]{beamer}
\newif\ifbibliography
\usepackage{pgfpages}
\setbeamertemplate{caption}[numbered]
\setbeamertemplate{caption label separator}{: }
\setbeamercolor{caption name}{fg=normal text.fg}
\beamertemplatenavigationsymbolsempty
% remove section numbering
\setbeamertemplate{part page}{
  \centering
  \begin{beamercolorbox}[sep=16pt,center]{part title}
    \usebeamerfont{part title}\insertpart\par
  \end{beamercolorbox}
}
\setbeamertemplate{section page}{
  \centering
  \begin{beamercolorbox}[sep=12pt,center]{section title}
    \usebeamerfont{section title}\insertsection\par
  \end{beamercolorbox}
}
\setbeamertemplate{subsection page}{
  \centering
  \begin{beamercolorbox}[sep=8pt,center]{subsection title}
    \usebeamerfont{subsection title}\insertsubsection\par
  \end{beamercolorbox}
}
% Prevent slide breaks in the middle of a paragraph
\widowpenalties 1 10000
\raggedbottom
\AtBeginPart{
  \frame{\partpage}
}
\AtBeginSection{
  \ifbibliography
  \else
    \frame{\sectionpage}
  \fi
}
\AtBeginSubsection{
  \frame{\subsectionpage}
}
\usepackage{iftex}
\ifPDFTeX
  \usepackage[T1]{fontenc}
  \usepackage[utf8]{inputenc}
  \usepackage{textcomp} % provide euro and other symbols
\else % if luatex or xetex
  \usepackage{unicode-math} % this also loads fontspec
  \defaultfontfeatures{Scale=MatchLowercase}
  \defaultfontfeatures[\rmfamily]{Ligatures=TeX,Scale=1}
\fi
\usepackage{lmodern}
\ifPDFTeX\else
  % xetex/luatex font selection
\fi
% Use upquote if available, for straight quotes in verbatim environments
\IfFileExists{upquote.sty}{\usepackage{upquote}}{}
\IfFileExists{microtype.sty}{% use microtype if available
  \usepackage[]{microtype}
  \UseMicrotypeSet[protrusion]{basicmath} % disable protrusion for tt fonts
}{}
\makeatletter
\@ifundefined{KOMAClassName}{% if non-KOMA class
  \IfFileExists{parskip.sty}{%
    \usepackage{parskip}
  }{% else
    \setlength{\parindent}{0pt}
    \setlength{\parskip}{6pt plus 2pt minus 1pt}}
}{% if KOMA class
  \KOMAoptions{parskip=half}}
\makeatother
\usepackage{longtable,booktabs,array}
\usepackage{caption}
\captionsetup[table]{skip=6pt}
\newcounter{none} % for unnumbered tables
\usepackage{calc} % for calculating minipage widths
\usepackage{caption}
% Make caption package work with longtable
\makeatletter
\def\fnum@table{\tablename~\thetable}
\makeatother
\setlength{\emergencystretch}{3em} % prevent overfull lines
\providecommand{\tightlist}{%
  \setlength{\itemsep}{0pt}\setlength{\parskip}{0pt}}
% Auto-generated by infrastructure/rendering/_pdf_latex_helpers.py::extract_math_font_preamble.
% Minimal header passed to Pandoc via -H so Beamer slides pick up the same math font
% as the combined-PDF path. Adjust the manuscript's preamble.md to change the font globally.
\usepackage{unicode-math}
\setmathfont{latinmodern-math.otf}

% Manuscript macro/package fallbacks (newcommand -> providecommand).
% Core mathematics
\usepackage{amsmath}
\usepackage{amssymb}
% Document layout
% Tables
\usepackage{booktabs}
% Caption styling (booktabs already loaded above). Small caption font with a
% bold label ("Figure 1", "Table 2") gives the math/figure-heavy layout a
% consistent, journal-like caption rhythm without fighting pandoc-crossref —
% pandoc emits standard \caption{}, and caption/captionsetup only restyle it.
% ── Theorem environments ─────────────────────────────────────────────────
% amsthm is loaded above. The FedGVI<->Friston bridge is stated as a small
% number of theorems/definitions/lemmas/propositions/corollaries; sharing one
% counter (definition/lemma/proposition/corollary all step the `theorem`
% counter) gives a single monotone numbering 1, 2, 3, ... across all kinds, so
% authors can reference them as Theorem 1, Definition 2, Corollary 3 without a
% per-kind counter drifting out of order. Plain style for theorem-like results,
% definition style (upright body) for definitions.
% (algorithm2e intentionally not loaded: this manuscript states results as
% numbered amsthm Theorems/Definitions, not algorithm2e pseudocode.)
% Code listings
\usepackage{listings}
% Typography and formatting
% (siunitx intentionally not loaded: no \SI/\num units appear in this manuscript.)
% Cross-references and citations
% (cleveref and titlesec intentionally not loaded: unavailable in this TeX tree
% and not required. Cross-references resolve through pandoc-crossref's own
% \ref-style names — no \cref appears — and section heading styling falls back
% to the pandoc/LaTeX defaults.)
% ── TeX Live Basic-compatible mono font for code listings ────────────
% Latin Modern Mono is bundled with TeX Live Basic and keeps the manuscript
% renderable on minimal TeX installations. Mathematical Unicode coverage is
% provided separately by Latin Modern Math below, so the code font does not
% need a private JuliaMono installation.
%
% Use the TeX font filename rather than the system-family name: the minimal
% TeX Live fontconfig database may not register the OpenType family even though
% kpsewhich can resolve the bundled files.
% Math font for unicode-math: Latin Modern Math (TeX Live) has full BMP
% coverage including U+2223 (\mid), U+226A/226B (\ll/\gg), and the Greek/
% blackboard letters used in equations. Without an explicit \setmathfont,
% unicode-math falls back to lmroman text font which lacks several glyphs
% and emits "Missing character" warnings on every \mid in math mode.

% Slide layout defaults for warning-clean scientific decks.
\usepackage{etoolbox}
\IfFileExists{xurl.sty}{\usepackage{xurl}}{}
\IfFileExists{seqsplit.sty}{\usepackage{seqsplit}}{\newcommand{\seqsplit}[1]{#1}}
\protected\def\breakseq#1{\seqsplit{#1}}
\protected\def\breaktt#1{\begingroup\ttfamily\seqsplit{#1}\endgroup}
\setlength{\emergencystretch}{6em}
\tolerance=5000
\hbadness=10000
\hfuzz=1pt
\setlength{\tabcolsep}{2pt}
\AtBeginEnvironment{longtable}{\tiny\renewcommand{\arraystretch}{0.86}\setlength{\tabcolsep}{1pt}}
\AtBeginEnvironment{tabular}{\tiny\renewcommand{\arraystretch}{0.86}\setlength{\tabcolsep}{1pt}}
\AtBeginEnvironment{equation}{\tiny}
\AtBeginEnvironment{equation*}{\tiny}
\AtBeginEnvironment{align}{\tiny}
\AtBeginEnvironment{align*}{\tiny}
\AtBeginEnvironment{itemize}{\footnotesize}
\AtBeginEnvironment{enumerate}{\footnotesize}
\AtBeginEnvironment{description}{\footnotesize}
\setbeamerfont{caption}{size=\tiny}
\setbeamerfont{caption name}{size=\tiny}
\setbeamerfont{normal text}{size=\small}
\setbeamerfont{frametitle}{size=\small}
\setbeamerfont{section title}{size=\footnotesize}
\setbeamerfont{subsection title}{size=\footnotesize}
\setbeamertemplate{section page}{%
  \centering
  \begin{beamercolorbox}[sep=12pt,center,wd=\paperwidth]{section title}
    \parbox{0.86\paperwidth}{\centering\usebeamerfont{section title}\insertsection\par}
  \end{beamercolorbox}
}
\setbeamertemplate{subsection page}{%
  \centering
  \begin{beamercolorbox}[sep=8pt,center,wd=\paperwidth]{subsection title}
    \parbox{0.86\paperwidth}{\centering\usebeamerfont{subsection title}\insertsubsection\par}
  \end{beamercolorbox}
}
\setlength{\abovecaptionskip}{2pt}
\setlength{\belowcaptionskip}{0pt}

% Natbib and cross-reference fallbacks — slides don't load natbib
% or cleveref, but combined-PDF manuscript prose may emit these
% commands. The fallback renders citations as a bracketed key list
% and cross-references as detokenized labels so slides don't fail on
% undefined control sequences or raw underscores. \providecommand is
% a no-op if packages load later.
\providecommand{\citep}[1]{[#1]}
\providecommand{\citet}[1]{#1}
\providecommand{\citealp}[1]{#1}
\providecommand{\citeauthor}[1]{#1}
\providecommand{\citeyear}[1]{#1}
\providecommand{\cref}[1]{\texttt{\detokenize{#1}}}
\providecommand{\Cref}[1]{\texttt{\detokenize{#1}}}
\providecommand{\crefrange}[2]{\texttt{\detokenize{#1}}--\texttt{\detokenize{#2}}}
\providecommand{\Crefrange}[2]{\texttt{\detokenize{#1}}--\texttt{\detokenize{#2}}}
\providecommand{\paragraph}[1]{\textbf{#1}\ }

% Opt-in accessible presentation profile.
\setbeamerfont{normal text}{size*={20pt}{24pt}}
\setbeamerfont{frametitle}{size*={28pt}{32pt}}
\setbeamerfont{section title}{size*={28pt}{32pt}}
\setbeamerfont{subsection title}{size*={28pt}{32pt}}
\setbeamerfont{caption}{size*={16pt}{19pt}}
\setbeamerfont{caption name}{size*={16pt}{19pt}}
\setbeamerfont{itemize/enumerate body}{size*={20pt}{24pt}}
\setbeamerfont{itemize/enumerate subbody}{size*={20pt}{24pt}}
\setbeamerfont{itemize/enumerate subsubbody}{size*={20pt}{24pt}}
\setbeamerfont{itemize item}{size*={20pt}{24pt}}
\setbeamerfont{itemize subitem}{size*={20pt}{24pt}}
\setbeamerfont{itemize subsubitem}{size*={20pt}{24pt}}
\setbeamerfont{description body}{size*={20pt}{24pt}}
\setbeamerfont{description item}{size*={20pt}{24pt}}
\setbeamerfont{quote}{size*={20pt}{24pt}}
\setbeamerfont{quotation}{size*={20pt}{24pt}}
\AtBeginEnvironment{longtable}{\fontsize{20pt}{24pt}\selectfont}
\AtBeginEnvironment{tabular}{\fontsize{20pt}{24pt}\selectfont}
\AtBeginEnvironment{equation}{\fontsize{20pt}{24pt}\selectfont}
\AtBeginEnvironment{equation*}{\fontsize{20pt}{24pt}\selectfont}
\AtBeginEnvironment{align}{\fontsize{20pt}{24pt}\selectfont}
\AtBeginEnvironment{align*}{\fontsize{20pt}{24pt}\selectfont}
\AtBeginEnvironment{itemize}{\fontsize{20pt}{24pt}\selectfont}
\AtBeginEnvironment{enumerate}{\fontsize{20pt}{24pt}\selectfont}
\AtBeginEnvironment{description}{\fontsize{20pt}{24pt}\selectfont}
\AtBeginDocument{\usebeamerfont{normal text}}
\newlength{\AFSlideTableWidth}
% The fixed footer reservation is calibrated with the semantic seven-line regular-body budget.
\setbeamertemplate{footline}{%
  \leavevmode\vbox{%
    \hbox{%
      \begin{beamercolorbox}[wd=\paperwidth,ht=16pt,dp=3pt,center]{author in head/foot}%
        \fontsize{16pt}{19pt}\selectfont Untagged PDF derivative \textbar\ \href{../web/index.html}{HTML reader}%
      \end{beamercolorbox}%
    }%
    \vskip6pt%
  }%
}
% Accessible footer reserved height: 25pt.

\newtheorem{proposition}{Proposition}
\newtheorem{hypothesis}{Hypothesis}
\newtheorem{remark}[theorem]{Remark}
\newtheorem{axiom}{Axiom}
\newtheorem{property}{Property}
\usepackage{bookmark}
\IfFileExists{xurl.sty}{\usepackage{xurl}}{} % add URL line breaks if available
\urlstyle{same}
% fallback for those not using the hyperref driver hyperxmp:
\makeatletter
\@ifundefined{xmpquote}{\newcommand{\xmpquote}[1]{#1}}{}
\makeatother
\hypersetup{
  hidelinks,
  pdfcreator={LaTeX via pandoc}}

\author{\texorpdfstring{}{}}
\date{}

\begin{document}

\section{Results: recovery checks and study suite}\label{sec:results}

\begin{frame}{Results: recovery checks and study suite}
Every quantitative assertion in this and the following results sections
is a generated token,
\end{frame}

\begin{frame}[fragile]{Results: recovery checks\ldots{} (part 2)}
hydrated by the manuscript-variable generator from analysis outputs
produced by \texttt{src/analysis/workflow.py} and
\texttt{src/fedference/experiments/} --- no number is transcribed by
hand.
\end{frame}

\begin{frame}{Results: recovery checks\ldots{} (part 3)}
The studies implement categorical source-mechanism analogues of the
colony belief-sharing scenario (Friston et al. 2024) and add the
contaminated-sentinel robustness sweep and the federated neural-network
baseline that are this paper's robust-federated-learning contribution.
\end{frame}

\begin{frame}{Results: recovery checks\ldots{} (part 4)}
All runs are deterministic under seed 0.
\end{frame}

\begin{frame}{Results: recovery checks\ldots{} (part 5)}
We lead with the recovery limits, not with a study, because they are
what makes the studies a single coherent system rather than a collection
of unrelated experiments. The generalized-Bayes machinery of
sec.~5 contains the standard-Bayes client corner,
\end{frame}

\begin{frame}{Results: recovery checks\ldots{} (part 6)}
while the server has the exact project-local zero-robustness
log-linear-pool identity.
\end{frame}

\begin{frame}{Results: recovery checks\ldots{} (part 7)}
Under the explicit bridge in sec.~5.8, that
pool is a categorical specialization of Eq. 7's message-combination term
rather than the complete source protocol. We verify those limited
identities to machine precision before reporting anything built on top
of them.
\end{frame}

\begin{frame}{Recovery limits: standard-Bayes and project-pool corners
are exact to machine precision}
\protect\phantomsection\label{sec:results-recovery}
\end{frame}

\begin{frame}{Recovery limits (part 2)}
The identities that anchor every result are the client recovery of
standard Bayes at the KL/NLL/\(\beta\to0\) and \(q_{\text{loss}}\to0\)
limits plus the project-local server recovery to the log-linear pool at
\(c=0\) --- the scoped claims of sec.~6 (Corollary
6,
\end{frame}

\begin{frame}{Recovery limits (part 3)}
Lemma 3, Theorem
5).

These are not figures but exact equalities, pinned by the locked core
test suite.
\end{frame}

\begin{frame}{Recovery limits (part 4)}
Under the theorem's shared-support, posterior-log-potential, and
fixed-weight assumptions, the server pool specializes Eq. 7's
message-combination term; it does not reproduce the source construction
in full. Robustness is a tested extension that vanishes at the stated
recovery limits.
\end{frame}

\begin{frame}{Recovery limits (part 5)}
The five residuals below are the maximum absolute deviations between
each generalized-Bayes object and the standard object it must reproduce
in the trusting limit.
\end{frame}

\begin{frame}{Recovery limits (part 6)}
Each is a deterministic constant of the mathematics, not a per-run
sample: it is reported as the maximum absolute deviation over the
recovery band, which is exactly \(0\) where the implementation evaluates
the closed form at the limit (the Rényi/loss switch) and otherwise a
tiny floating-point residual:
\end{frame}

\begin{frame}{Recovery limits (part 7)}
\textbf{Server recovery.} The server-side aggregator at zero robustness
equals the log-linear pool (eq.~10, Theorem
5): maximum absolute deviation 0. This
is the \emph{naive-recovery} limit of the server-side heuristic ---
\end{frame}

\begin{frame}{Recovery limits (part 8)}
the only property proven for that axis (see
sec.~3.3).
\end{frame}

\begin{frame}{Recovery limits (part 9)}
\textbf{Posterior recovery.} The generalized posterior under the KL
divergence and the NLL loss equals the closed-form
prior\(\times\)likelihood Bayes posterior (eq.~11,
Corollary 6): maximum absolute deviation
5.55e-17.
\end{frame}

\begin{frame}{Recovery limits (part 10)}
\textbf{Divergence recovery.} The Rényi divergence recovers KL as
\(\alpha\to1\) (eq.~6, Lemma
3): residual 0.
\end{frame}

\begin{frame}{Recovery limits (part 11)}
\textbf{Loss recoveries.} The \(\beta\)-loss recovers the NLL as
\(\beta\to 0\) (eq.~7, Proposition
4): residual 0. The robust categorical
cross-entropy recovers the NLL as \(q_{\text{loss}}\to 0\)
(eq.~8,
\end{frame}

\begin{frame}{Recovery limits (part 12)}
Proposition 4): residual 0.
\end{frame}

\begin{frame}{Recovery limits (part 13)}
Because the Rényi divergence and the two categorical losses switch to
their exact closed form inside narrow numerical-stability bands around
the limit point (the Rényi switch band for \(\alpha\) and the
categorical-loss switch band for \(q_{\text{loss}}\) and \(\beta\)),
\end{frame}

\begin{frame}{Recovery limits (part 14)}
the three zero residuals above confirm that branch equals the standard
object --- not, by themselves, that the \emph{general} formula converges
there.
\end{frame}

\begin{frame}{Recovery limits (part 15)}
As a genuine (non-branch) convergence witness, evaluating each general
formula strictly \emph{outside} its switch band ---
\(q_{\text{loss}} = \beta =
1.00 \times 10^{-6}\) for the categorical losses and
\(\alpha = 1.00001\) for the Rényi divergence --- gives residuals
1.12e-05 (rcce), 1.24e-05 (\(\beta\)-loss), and 1.66e-05 (Rényi).
\end{frame}

\begin{frame}{Recovery limits (part 16)}
These residuals are nonzero---a small multiple of the input offset
itself, as the first-order Taylor behavior near the limit predicts---yet
remain several orders of magnitude below the \(O(1)\) scale of the
loss/divergence values being compared, and shrinking monotonically as
the offset shrinks toward the switch band.
\end{frame}

\begin{frame}[fragile]{Recovery limits (part 17)}
The verification is in \texttt{test\_core\_identities.py} under
\texttt{tests/fedference/}.

This is evidence that the general formula itself converges to the
standard-Bayes limit, not merely that the implementation switches to it
exactly at the corner.
\end{frame}

\begin{frame}{Recovery limits (part 18)}
The first residual is the naive-aggregate limit of the
\emph{server-side} heuristic (Theorem
5);
\end{frame}

\begin{frame}{Recovery limits (part 19)}
the latter four are the per-agent generalized-Bayes recoveries
(Corollary 6 + Proposition
4) and the divergence-family recovery in
the Rényi limit (Lemma 3,
\end{frame}

\begin{frame}{Recovery limits (part 20)}
eq.~6) that define the theorem-bearing FedGVI axis
under matching assumptions.
\end{frame}

\begin{frame}{Recovery limits (part 21)}
Keeping the three axes distinct at the level of the recovery limits is
what lets the robustness claims of sec.~7.5 and
sec.~7.6 rest on the per-agent axis without
leaning on the heuristic.
\end{frame}

\begin{frame}[fragile]{Recovery limits (part 22)}
266 of 268 acceptance criteria are verified. The pure-NumPy/SciPy core
carries project test coverage of 90.57\% (gate \(\ge 90\%\)), with every
stochastic step threaded through a single seeded
\texttt{np.random.default\_rng(0)}. sec.~10
records the full environment fingerprint,
\end{frame}

\begin{frame}{Recovery limits (part 23)}
and the expected-free-energy identity that underwrites the
active-inference substrate is proven and visualized in
sec.~6.2 (fig.~8).
\end{frame}

\end{document}
